\documentclass{scrreprt}

\usepackage[ngerman]{babel}		% Sprachen
\usepackage[latin1]{inputenc} % Eingabe von �,�,�,� erlaubt 
\usepackage[T1]{fontenc}			% => Umlauten koennen richtig getrennt werden

\usepackage{graphicx} % zum Laden von Bildern 
\begin{document}

\tableofcontents

\chapter{Kapitel 1}
\label{sec:Kapitel1}

\section{�berschrift 1.1}
%\label{sec:�berschrift11}
Text 1.1 Weit hinten, hinter den Wortbergen, fern der L�nder Vokalien und Konsonantien leben die Blindtexte. Abgeschieden wohnen Sie in Buchstabhausen an der K�ste des Semantik, eines gro�en Sprachozeans. Ein kleines B�chlein namens Duden flie�t durch ihren Ort und versorgt sie mit den n�tigen Regelialien. 
\begin{quote}
	Ein kurzes Zitat
\end{quote}
Es ist ein paradiesmatisches Land, in dem einem gebratene Satzteile in den Mund fliegen. Nicht einmal von der allm�chtigen Interpunktion werden die Blindtexte beherrscht - ein geradezu unorthographisches Leben. Eines Tages aber beschloss eine kleine Zeile Blindtext, ihr Name war Lorem Ipsum, hinaus zu gehen in die weite Grammatik.

\begin{table}[htbp]
	\centering
		\begin{tabular}{l|l l}
			11 & 12 & 13\\
			\hline
			21 & 22 & 23\\
			31 & 32 & 33\\
		\end{tabular}
	\caption{test-tabelle}
	\label{tab:testtabelle}
\end{table}

\subsection{Unter�berschrift 1.1.1}
%\label{sec:�berschrift111}
Text 1.1.1. Der gro�e Oxmox riet ihr davon ab, da es dort wimmele von b�sen Kommata, wilden Fragezeichen und hinterh�ltigen Semikoli, doch das Blindtextchen lie� sich nicht beirren. Es packte seine sieben Versalien, schob sich sein Initial in den G�rtel und machte sich auf den Weg. 
	\[E = m \cdot c^{2}\]
Als es die ersten H�gel des Kursivgebirges erklommen hatte, warf es einen letzten Blick zur�ck auf die Skyline seiner Heimatstadt Buchstabhausen, die Headline von Alphabetdorf und die Subline seiner eigenen Stra�e, der Zeilengasse. Wehm�tig lief ihm eine rhetorische Frage �ber die Wange, dann setzte es seinen Weg fort.

\begin{figure}[htbp]
	\centering
			\includegraphics[scale=0.50]{images/tu_logo}
	\caption{tu-logo}
	\label{fig:tulogo}
\end{figure}

\section{�berschrift 2}
%\label{sec:�berschrift2}

\begin{quotation}
	Eine langes abgesetztes Zitat: Unterwegs traf es eine Copy. Die Copy warnte das Blindtextchen, da, wo sie herk�me w�re sie zigmal umgeschrieben worden und alles, was von ihrem Ursprung noch �brig w�re, sei das Wort "und" und das Blindtextchen solle umkehren und wieder in sein eigenes, sicheres Land zur�ckkehren. Doch alles Gutzureden konnte es nicht �berzeugen und so dauerte es nicht lange, bis ihm ein paar heimt�ckische Werbetexter auflauerten, es mit Longe und Parole betrunken machten und es dann in ihre Agentur schleppten, wo sie es f�r ihre Projekte wieder und wieder missbrauchten. Und wenn es nicht umgeschrieben wurde, dann benutzen Sie es immer noch.
\end{quotation}

\section{�berschrift 3}
%\label{sec:�berschrift3}
Text 3 Weit hinten, hinter den Wortbergen, fern der L�nder Vokalien und Konsonantien leben die Blindtexte. Abgeschieden wohnen Sie in Buchstabhausen an der K�ste des Semantik, eines gro�en Sprachozeans. Ein kleines B�chlein namens Duden flie�t durch ihren Ort und versorgt sie mit den n�tigen Regelialien. 
	\[ m = \int^{1}_{x=0}{\frac{x + 3 \Pi}{\sqrt{\Phi}}}dx\]
Es ist ein paradiesmatisches Land, in dem einem gebratene Satzteile in den Mund fliegen. Nicht einmal von der allm�chtigen Interpunktion werden die Blindtexte beherrscht - ein geradezu unorthographisches Leben. Eines Tages aber beschloss eine kleine Zeile Blindtext, ihr Name war Lorem Ipsum, hinaus zu gehen in die weite Grammatik.

\subsection{Unterabschnitt}
%\label{sec:Unterabschnitt}
Eine Aufz�hlung:
\begin{itemize}
	\item erster Punkt
	\item zweiter Punkt
\end{itemize}
und eine Nummerierung:
 \begin{enumerate}
	\item erster Punkt
	\item zweiter Punkt
	\item dritter Punkt
\end{enumerate}

\subsubsection{Unter-Unterabschnitt}
%\label{sec:UnterUnterabschnitt}
Einen Begriffsdefinition: 
\begin{description}
	\item[Begriff] Ich bin eine l�ngere Beschreibung.
	\item[Begriff2] Ich auch.
\end{description}

\paragraph{Ein Absatz}
%\label{sec:Absatz}
 mit einem Gedicht
\begin{verse}
	Wer reitert so sp�t durch Nacht und Wind?\\
	Es ist der Vater mit seinem Kind.\\
\end{verse}
\subparagraph{Unterabsatz}
%\label{sec:Unterabsatz}
Text Text

\end{document} 
